\section{Scope and limitations of the assessment}
\label{sec:scope}

The experiment isolates explainer randomness for fixed trained classifiers. A fixed split and one model-training seed do not measure variability from retraining, hyperparameter selection or distribution shift. Three tabular tasks and the selected model families cannot establish a universal ranking of SHAP implementations. In particular, deterministic exact algorithms and newer stochastic estimators have different properties from the two approximations compared here. The logistic classifier is explained on its probability scale, so its linear predictor does not make these probability-game attributions analytically linear.

Exact enumeration removes coalition-sampling error from the selected finite-background reference, subject to floating-point accuracy and validated implementation. It does not make that empirical background representative or resolve the choice between marginal and observational games. Mixing feature blocks can produce unlikely records. A separately changed background defines a different target; its shift must be distinguished from estimation error for either target. The cost of enumeration grows exponentially with the number of original-feature players. Its feasibility for these modest-dimensional tasks does not supply a general reference strategy for high-dimensional deployment.

The selected magnitude floor and signed top-three criterion encode a numerical reading of an attribution, not a validated human decision rule. Other tolerances can change the adequacy rate, and small reference magnitudes or ties can make a feature list ambiguous. An input-panel mean does not certify each prediction, and thirty seeds cannot rule out rare failures for one input. Pointwise cluster intervals describe sampled-input variability under the stated independence assumption; they are not simultaneous bounds over every comparison. The conservative bounded-mean guard further limits any population reliability certificate, even when no failures are observed.

The data also constrain interpretation. The ACS benchmark is unweighted and samples person rows; original household or survey-design dependence is not modeled by the input bootstrap. The historical Diabetes data concern a narrow demographic sample. No result of this experiment establishes population fairness, causal validity, clinical performance or lending-policy suitability. More reproducible numerical attributions can still give a misleading rationale for a biased or poorly specified model. Reporting the numerical endpoint as a deployment endorsement would exceed the study's evidence.
